\section{Evaluation}

We keep the evaluation narrow and tied to the released implementation. Instead of relying on a single smoke-style benchmark, we report a fixed-grid experiment driven by the released Ghostext evaluation script \nolinkurl{run_prompt_message_grid.py}. The grid crosses ten public prompts with ten public messages for 100 end-to-end encode/decode trials under the same pinned local backend from \S\ref{sec:open-science}. The prompt set contains five English and five Chinese templates. The message set contains four short, four medium, and two long plaintexts spanning both languages. All trials use the released defaults. We keep seed $7$ throughout. The decoding settings remain \texttt{top\_p}=0.995, \texttt{max\_candidates}=64, \texttt{min\_entropy\_bits}=0.0, and \texttt{total\_frequency}=4096. The built-in low-entropy retry logic stays enabled throughout. A fixed local passphrase is reused across the grid for reproducibility, but only the policy and not the secret itself is released.

\subsection{Fixed 10x10 Prompt-Message Grid}

\begin{table}[t]
\caption{Fixed 10x10 prompt-message grid on the pinned local \texttt{llama-cpp} backend.}
\label{tab:grid-benchmark}
\footnotesize
\centering
\setlength{\tabcolsep}{2pt}
\begin{tabular}{@{}lrrrrrr@{}}
\toprule
Slice & Trials & Success & Enc med (s) & Dec med (s) & Mean bpt & Attempts mean \\
\midrule
Overall & 100 & 100/100 & 37.74 & 34.47 & 2.619 & 1.40 \\
EN prompts & 50 & 50/50 & 35.04 & 34.21 & 2.686 & 1.18 \\
ZH prompts & 50 & 50/50 & 38.90 & 37.09 & 2.552 & 1.62 \\
Short msgs & 40 & 40/40 & 27.90 & 26.45 & 2.645 & 1.08 \\
Medium msgs & 40 & 40/40 & 42.31 & 37.43 & 2.594 & 1.48 \\
Long msgs & 20 & 20/20 & 60.22 & 58.11 & 2.615 & 1.90 \\
\bottomrule
\end{tabular}
\end{table}

Table~\ref{tab:grid-benchmark} shows exact round-trip recovery for all 100 trials, with decode success rate $1.000$ and Wilson 95\% confidence interval $[0.963, 1.000]$. No decode mismatch or fail-open behavior appears in the released run log. Overall median encode latency is $37.74$\,s and median decode latency is $34.47$\,s. Mean encode payload density is $2.619$ bits/token (median $2.677$), while mean packet length is $167.1$ bytes.

The main performance caveat is not correctness but tail latency. Encode latency has mean $50.56$\,s, $p90=74.69$\,s, $p99=196.55$\,s, and maximum $291.69$\,s, whereas decode latency is tighter with mean $38.02$\,s, $p90=58.07$\,s, and $p99=72.22$\,s. Retry behavior explains much of that spread, because some local contexts trigger low-entropy or no-stable-candidate retries. English-prompt cases are modestly faster and denser than Chinese-prompt cases, while message length mainly affects latency rather than mean bits/token.

These results support one narrow claim: within one pinned local environment and one fixed prompt-message grid, Ghostext provides deterministic, fail-closed end-to-end recoverability with transparent overhead measurements. They also give a more realistic view of runtime than the earlier smoke benchmark, because they expose retry-driven tail behavior across multiple prompt and payload regimes. They do not establish detector resistance, direct distribution-distance bounds, cross-hardware determinism, or robustness to active text modification.
